% !TeX root = ../../../navier-stokes-manuscript.tex
% ============================================================
% 1.1 The Regularity Problem
% ============================================================
\subsection{The Regularity Problem}\label{sub:TheRegularityProblem}

The three-dimensional Navier--Stokes regularity problem asks whether a smooth solution can lose its smoothness in finite time.
Fix smooth, divergence-free, finite-energy initial data \(u_0\), fix \(\nu>0\), and let \((u,p)\) be the corresponding maximal classical solution on \([0,T_*)\).
On \(\Omega=\mathbb R^3\) or \(\mathbb T^3\), the pair satisfies
\[
  \partial_tu+(u\cdot\nabla)u+\nabla p
  =
  \nu\Delta u,
  \qquad
  \nabla\cdot u=0,
  \qquad
  u(x,0)=u_0(x).
\]
We want to prove that \(T_*\) cannot be finite.

In the whole-space formulation, the datum and its spatial derivatives decay rapidly at infinity.
In the periodic formulation, \(u_0\), \(u\), and \(p\) are periodic in the spatial variables.
These are the positive existence-and-smoothness alternatives in the Clay problem \parencite{Fefferman2000}.

\divider{Smoothness and Maximal Time}

\begin{definition}[Smoothness through a finite time]\label{def:Smoothness}
The velocity--pressure pair is smooth through \(T>0\) when
\[
  u\in C^\infty\!\left(\Omega\times[0,T];\mathbb R^3\right),
  \qquad
  p\in C^\infty\!\left(\Omega\times[0,T]\right).
\]
Equivalently, every mixed spatial and temporal derivative of each finite order exists and is continuous through \(t=T\).
\end{definition}

\begin{definition}[Maximal smooth time]\label{def:Continuation}
The maximal smooth time \(T_*\) is the supremum of the times \(T\) for which the local solution generated by \(u_0\) extends as a smooth solution on \([0,T)\).
\end{definition}

If \(T_*<\infty\), there is no \(\varepsilon>0\) for which the same classical solution extends smoothly to \([0,T_*+\varepsilon)\).
Global regularity is exactly the statement \(T_*=\infty\) for every admissible datum.
